\documentclass[11pt,a4paper]{article}
\usepackage[a4paper,margin=2.3cm]{geometry}
\usepackage{fontspec}
\usepackage[spanish]{babel}
\usepackage{amsmath,amssymb}
\usepackage{xcolor}
\usepackage{enumitem}
\usepackage{microtype}
\usepackage[hidelinks]{hyperref}
\setlength{\parindent}{0pt}
\setlength{\parskip}{4pt}
\definecolor{azul}{RGB}{30,58,95}
\definecolor{gris}{RGB}{90,90,90}

\newcommand{\fuente}[1]{\par{\small\color{gris}Fuente: #1}}
\newcommand{\seccion}[1]{\vspace{0.4cm}{\Large\color{azul}\textbf{#1}}\par\vspace{0.1cm}}
\newcommand{\bem}{$\flat$}
\setlist[enumerate]{leftmargin=*,itemsep=6pt}
\setlist[itemize]{itemsep=1pt,topsep=2pt}

\begin{document}

{\LARGE\color{azul}\textbf{Puntos para la defensa}}\par
\vspace{0.1cm}
{\large Análisis computacional microtonal de Angine de Poitrine}\par
\vspace{0.3cm}
Ordenados como avanza una defensa: la herramienta, la armonía, la melodía, la forma, la síntesis y las preguntas
difíciles. Cada punto dice de dónde sale la cifra para poder comprobarla. Las gráficas 1a--3b son las de
\texttt{make graficas-diapositivas}; los § remiten a \texttt{docs/evolucion\_de\_la\_tesis.md}.

\seccion{1. La herramienta: por qué creerle}
\begin{enumerate}
\item \textbf{Una herramienta que no existía.} Ninguna herramienta estándar transcribe en 24-EDO (cuartos de tono).
  El cuaderno separa guitarra y bajo, sigue la altura, detecta compases y ciclos del looper y analiza la armonía.
  \fuente{\texttt{docs/evolucion\_de\_la\_tesis.md} §1.}
\item \textbf{Se mide contra partituras corregidas de oído, no contra la oficial.} La de Guitar Pro también tiene
  errores, y la de Fabienk es de la versión en vivo.
  \fuente{§2.}
\item \textbf{Prueba a ciegas.} Yor Zarad se reservó sin abrir ninguna salida del cuaderno: 36.8 en la melodía,
  47.7 sin contar octavas y 68.6 en el bajo; el centro de la canción queda al nivel de las demás.
  \fuente{§4.}
\item \textbf{Arbitraje a ciegas.} En la primera ronda la partitura tenía razón en la mitad de los desacuerdos.
  Desde entonces, que el cuaderno y la partitura no coincidan no basta para corregir.
  \fuente{§2.}
\item \textbf{La herramienta no usa la partitura para guiarse.} Se quitó todo lo que dependía del .gp, así que sirve
  para canciones sin partitura.
  \fuente{§2.}
\item \textbf{Reproducible.} Demucs usa semilla fija y se guardan huellas SHA-256 del audio. Un hallazgo solo cuenta
  si se mantiene con dos separaciones distintas (la misma canción dio 57 y 62).
  \fuente{§2.}
\item \textbf{Un separador mejor no cambió el resultado.} BS-RoFormer tiene mejor SDR público que Demucs (9.05 frente
  a 5.22 en guitarra), pero en Sarniezz dio el mismo 44.9 en la melodía y tarda una hora por canción. El cuello de
  botella no es la separación sino las capas superpuestas del looper. Responde de antemano a «¿por qué no usaste un
  modelo mejor?».
  \fuente{\texttt{docs/separacion\_roformer.md}.}
\end{enumerate}

\seccion{2. Armonía: no encaja en ningún estilo medido}
\begin{enumerate}[resume]
\item \textbf{Ninguna canción de la banda es típica de ningún corpus.} El máximo es 17.7 (Utzp en pop); Fabienk,
  Yor Zarad y Sarniezz no pasan de 7.2.
  \fuente{matriz de verosimilitud, gráfica 1a.}
\item \textbf{El contraste da peso al resultado.} Con el mismo método, las referencias sí caen en estilos conocidos:
  Gentle Giant «Experience» 92.4 en pop, King Gizzard «Billabong Valley» 86.2 en rock, Primus «Tommy the Cat» 75.1
  en pop. El método no declara todo «raro».
  \fuente{matriz de las 36 canciones.}
\item \textbf{El método está calibrado.} Sobre 388 piezas reales de rock, las clasifica como «estilo existente» en el
  92.3\,\% de los casos y como «sin vecino» solo en el 5.2\,\%. En la banda, «sin vecino» aparece en 9 de 24 casos
  (37.5\,\%), unas siete veces más. Probablemente el mejor argumento contra «tu método siempre dice que es nuevo».
  \fuente{salida de \texttt{verosimilitud\_por\_bloques.py}.}
\item \textbf{Sarniezz y Fabienk caen donde no hay piezas de ningún corpus,} por el peso de \bem VII y \bem III (parte
  baja del mapa). Mata Zyklek y Utzp, que son pedales, quedan en el borde del rock.
  \fuente{mapa de fundamentales, gráfica 1b, y §6.}
\item \textbf{Distancia al pop de Hooktheory:} 0.40--0.59 para la banda; los álbumes de pop entre sí no pasan de 0.07.
  \fuente{§6.}
\item \textbf{Los umbrales ad hoc se descartaron.} Las primeras pruebas de estilo usaban umbrales propios y se
  reemplazaron por corpus reales. Sirvieron para descartar jazz, práctica común, tonos enteros y dodecafonismo.
  \fuente{§6.}
\end{enumerate}

\seccion{3. Melodía: el elemento nuevo}
\begin{enumerate}[resume]
\item \textbf{Los cuartos de tono son grados de la escala, no adornos.} La banda concentra entre el 48 y el 82\,\% de
  sus cuartos de tono en dos grados fijos, que vuelven en casi todas las vueltas del ciclo.
  \fuente{gráfica 2a y §7.}
\item \textbf{Piso de ruido medido.} En Primus, que no usa cuartos de tono, el cuaderno marca el 9--13\,\% de la
  melodía en cuartos de tono por bends y slides. Lo que supera ese piso, como grado, es de la banda.
  \fuente{gráfica 2a.}
\item \textbf{Dos familias de escala.}
  \begin{itemize}
  \item \bem 2+ y \bem 6+ sobre un modo frigio: Mata Zyklek, Yor Zarad y Utzp. Cercana al maqam y al makam turco, que
    la banda cita.
  \item 7+, 6+ y \bem 3+ sobre un modo dórico: Sarniezz y Fabienk.
  \end{itemize}
  \fuente{círculo de 24-EDO, gráfica 2b, y §7.}
\item \textbf{Base matemática del círculo de 24-EDO.} La quinta mide 14 pasos y $\gcd(14,24)=2$: el círculo se parte
  en dos campos, los 12 grados regulares y los 12 en cuarto de tono. En 19, 31 o 53-EDO hay un solo campo. Une el
  álgebra del artículo con la música.
  \fuente{gráfica 2b y la sección «Estructura algebraica» del artículo.}
\item \textbf{Angor es la excepción:} solo un 5\,\% de cuartos de tono, bajo el piso de ruido. Si preguntan por qué
  una canción no sigue el patrón, la respuesta está en los datos.
  \fuente{gráfica 2a.}
\end{enumerate}

\seccion{4. Forma: los ciclos del looper}
\begin{enumerate}[resume]
\item \textbf{Picos altos y aislados en la repetición.} La banda tiene picos en 20, 24, 28 y 32 corcheas; Primus y
  Zappa tienen curvas planas.
  \fuente{gráfica 3a.}
\item \textbf{Los picos coinciden con lo que dice la banda (hipótesis).} En el podcast hablan de tocar en 5 y 7
  tiempos:
  \begin{itemize}
  \item Fabienk tiene casi todos sus compases de 7 pulsos; su pico de 28 corcheas equivale a dos compases de 7.
  \item Yor Zarad tiene tramos de 7 y el mismo pico de 28.
  \item Mata Zyklek tiene tramos de 5; su pico de 20 corcheas equivale a dos compases de 5, con picos menores en 5 y 10.
  \end{itemize}
  El cuaderno llegó a eso desde el audio, sin que se lo dijeran.
  \fuente{\texttt{ciclos\_metricos.md} y §11.}
\item \textbf{Sarniezz y Utzp, en cambio, son de 4:} picos en 16 y 32 y tramos de 4 y 8 pulsos. No toda la banda es
  de compases irregulares; cerca de la mitad, como dice la tesis.
  \fuente{\texttt{ciclos\_metricos.md}.}
\end{enumerate}

\seccion{5. Síntesis e influencias}
\begin{enumerate}[resume]
\item \textbf{El híbrido se ve bloque por bloque.} La armonía y el metro apuntan hacia rock; la escala y los cuartos
  de tono, hacia radif, cada uno con un pariente distinto. Rock gana 8 de 24 bloques, radif 7 y en 9 no hay vecino.
  \fuente{gráfica 3b.}
\item \textbf{El vecino más cercano es King Gizzard, que la banda cita como influencia.} En armonía, la canción más
  parecida de cada tema de la banda es casi siempre de King Gizzard:
  \begin{itemize}
  \item Angor con «Flying Microtonal Banana»: distancia 0.00, idénticas (las dos son un pedal 100\,\% sobre I; caen en
    el mismo punto del mapa y tienen los mismos percentiles);
  \item Utzp con la misma canción: 0.01;
  \item Mata Zyklek con «Open Water»: 0.03.
  \end{itemize}
  En metro pasa lo mismo: Utzp, Yor Zarad y Sarniezz están a 0.00--0.01 de canciones de King Gizzard. El análisis
  confirma una influencia declarada.
  \fuente{salida de \texttt{verosimilitud\_por\_bloques.py}.}
\item \textbf{Sarniezz y Fabienk no tienen vecino ni siquiera en King Gizzard.} En armonía, su canción más parecida
  queda a 0.47--0.51. Lo propio de la banda está ahí.
  \fuente{la misma salida.}
\item \textbf{Primus es una comparación propia, no una influencia declarada.} Ninguna fuente pública la cita. Comparte
  con Fabienk y Sarniezz el centro, el modo dórico y las fundamentales 1, \bem 3 y \bem 7, pero en todos los bloques
  queda más lejos que el vecino más cercano. Conviene presentarlo así antes de que lo señalen.
  \fuente{§8.}
\end{enumerate}

\seccion{6. Preguntas difíciles}
\begin{enumerate}[resume]
\item \textbf{«King Gizzard también concentra los cuartos de tono.»} Es cierto: Melting 96\,\%, Doom City 95\,\%,
  Sleep Drifter 99.5\,\%. Lo que distingue a la banda no es concentrarlos, sino qué grados fija (\bem 2+ y \bem 6+, o
  7+ y 6+) y sobre qué modo. Además, las canciones microtonales de King Gizzard también salen atípicas en la matriz
  (Flying Microtonal Banana 9.1, Sleep Drifter 7.9): la banda pertenece a una corriente microtonal en el rock, no está
  sola.
\item \textbf{«¿Cuánto acierta la transcripción?»} Melodía: 66.4 en Mata Zyklek, 46.2 en Fabienk, 44.9 en Sarniezz y
  36.8 en la prueba a ciegas. Bajo: 69--83. Fundamental correcta: 95\,\% en Mata Zyklek, 82--83\,\% en Fabienk y
  58\,\% en Sarniezz (el punto débil, por su riff cromático en cuartos de tono). Los resultados de estilo usan sobre
  todo la armonía y los ciclos, que son lo más confiable.
\item \textbf{«Los corpus no tienen cuartos de tono.»} Es cierto. Las fundamentales en cuarto de tono se llevan al
  grado de abajo y su peso se informa aparte (el «¼\,\%» de la matriz: 14.4\,\% en Yor Zarad, 3.9\,\% en Sarniezz).
  Si eso sesga algo, es hacia parecerse más a los corpus: la conclusión de que no encaja es conservadora.
\item \textbf{«Son pocas canciones.»} Son seis, todo el Vol.~II, y 30 de referencia medidas con el mismo cuaderno.
\item \textbf{Límites que conviene decir primero:}
  \begin{itemize}
  \item fuera de la banda falla la elección del centro (Tommy the Cat sale en Sol en vez de Mi);
  \item la síncopa del radar rítmico da lo mismo que el azar;
  \item la referencia de Utzp no pasó la verificación de oído.
  \end{itemize}
  \fuente{§5, §9 y §10.}
\end{enumerate}

\vspace{0.3cm}
\colorbox{azul!8}{\parbox{\dimexpr\linewidth-2\fboxsep}{%
\textbf{El hilo en una frase.} La banda toma la forma del rock en loop, la armonía del rock modal y el metro de King
Gizzard, y le suma grados neutros fijos en dos familias: la frigia (\bem 2+, \bem 6+), que comparte con el radif persa
y con King Gizzard, y la dórica (7+, 6+, \bem 3+), que ningún corpus medido tiene. El método no
exagera la novedad: en rock real solo declara «sin vecino» el 5\,\% de las veces.}}

\end{document}
